\documentclass[11pt,a4paper]{article}
\usepackage[margin=2cm]{geometry}
\usepackage{booktabs}
\usepackage{xcolor}
\usepackage{colortbl}
\usepackage{array}
\usepackage{enumitem}
\usepackage{parskip}
\usepackage{titlesec}
\usepackage{mdframed}
\usepackage{microtype}

\definecolor{hdr}{RGB}{30,70,130}
\definecolor{boxbg}{RGB}{240,248,255}
\definecolor{wherebg}{RGB}{235,255,235}
\definecolor{rowgray}{RGB}{245,245,245}

\newcolumntype{L}[1]{>{\raggedright\arraybackslash}p{#1}}

\titleformat{\section}{\large\bfseries\color{hdr}}{}{0em}{}[\vspace{3pt}\hrule]
\titlespacing*{\section}{0pt}{14pt}{8pt}

\newmdenv[backgroundcolor=boxbg, linecolor=hdr, linewidth=1pt,
          leftmargin=0pt, rightmargin=0pt, innerleftmargin=10pt,
          innerrightmargin=10pt, innertopmargin=8pt, innerbottommargin=8pt,
          skipabove=6pt, skipbelow=4pt]{conceptbox}

\newmdenv[backgroundcolor=wherebg, linecolor=green!50!black, linewidth=1pt,
          leftmargin=0pt, rightmargin=0pt, innerleftmargin=10pt,
          innerrightmargin=10pt, innertopmargin=6pt, innerbottommargin=6pt,
          skipabove=4pt, skipbelow=6pt]{wherebox}

\setlist[itemize]{noitemsep, topsep=2pt, leftmargin=1.4em}

\begin{document}

\begin{center}
{\Large\bfseries\color{hdr} Statistical Tests Used in the Thesis}\\[4pt]
{\normalsize Explained simply, with their exact location in the analysis}
\end{center}
\vspace{-4pt}\hrule\vspace{10pt}

{\small This document explains every statistical test used in the thesis pipeline in plain language.
No prior statistics knowledge is assumed. Each test has a \textbf{What it is} block (the concept),
a \textbf{Simple example} (a real-world analogy), and a \textbf{Where we use it} block
(exactly which script and research question).}

\vspace{6pt}

%% ═══════════════════════════════════════════════════════════════════════════
\section{Wilson 95\% Confidence Interval (CI)}
%% ═══════════════════════════════════════════════════════════════════════════

\begin{conceptbox}
\textbf{What it is.}
When you calculate an adoption rate (e.g.\ ``30\% of repos in France are agentic''), that number
comes from a limited sample. The Wilson CI gives you a \emph{range} around that number and says:
``we are 95\% sure the true rate falls somewhere between X\% and Y\%''. The smaller the sample,
the wider the range. The larger the sample, the tighter and more trustworthy the range.

\medskip
\textbf{Simple example.}
You ask 10 people if they like pizza and 7 say yes --- that is 70\%. But with only 10 people,
you are not very confident. The CI might say: ``somewhere between 35\% and 93\%''. Now ask
1000 people and get 70\% --- the CI becomes ``67\% to 73\%''. Same percentage, much more trust.

\medskip
\textbf{Why Wilson specifically?}
The standard formula for CIs breaks down when percentages are very close to 0\% or 100\%,
or when samples are small. The Wilson formula is more accurate in those edge cases.
\end{conceptbox}

\begin{wherebox}
\textbf{Where we use it.}
\begin{itemize}
  \item \textbf{Stage 4a (Geographic analysis):} every country gets an adoption rate + Wilson CI.
        Countries with fewer than 30 repos are excluded because their CI would be too wide to be meaningful.
  \item \textbf{Stage 5g (World regions):} each of the 10 macro-regions (North America, Europe,
        East Asia, etc.) gets an adoption rate + Wilson CI.
  \item \textbf{Stage 4b (Top/Bottom countries):} the dot-plot of top and bottom countries
        shows the Wilson CI as error bars around each country's rate.
\end{itemize}
\end{wherebox}

%% ═══════════════════════════════════════════════════════════════════════════
\section{Binomial Test}
%% ═══════════════════════════════════════════════════════════════════════════

\begin{conceptbox}
\textbf{What it is.}
A binomial test checks whether an observed percentage is significantly higher (or lower) than
a specific baseline you choose in advance. You set a null hypothesis --- a boring default
assumption --- and ask: ``is my result too high to be a coincidence?''

\medskip
\textbf{Simple example.}
You claim a coin is fair (50\% heads). You flip it 100 times and get 68 heads. The binomial
test tells you: if the coin were truly fair, would you get 68 or more heads by pure luck?
If the answer is ``less than 5\% chance'', you reject the fair-coin claim.

\medskip
In the thesis, the baseline is 5\%. The question is: ``is global AI adoption significantly
higher than a trivial 5\%?'' The test confirms this with a p-value far below 0.05.
\end{conceptbox}

\begin{wherebox}
\textbf{Where we use it.}
\begin{itemize}
  \item \textbf{Stage 4b, RQ1 --- Global adoption:} tests whether the observed global adoption
        rate is significantly greater than the 5\% null hypothesis baseline.
        Reports the rate, Wilson CI, and p-value.
\end{itemize}
\end{wherebox}

%% ═══════════════════════════════════════════════════════════════════════════
\section{McNemar Test}
%% ═══════════════════════════════════════════════════════════════════════════

\begin{conceptbox}
\textbf{What it is.}
McNemar is used when you apply two different methods to the \emph{same} set of items and want
to know if the two methods give meaningfully different results. Because both methods look at
the same repos, they are not independent --- McNemar accounts for this pairing.

\medskip
\textbf{Simple example.}
A doctor tests 200 patients for a disease using two different tests (Test A and Test B) on the
same patients. Some patients are positive on A but negative on B, and vice versa. McNemar
checks whether the two tests disagree \emph{more than you would expect by chance}.

\medskip
In the thesis, the two ``tests'' are: (1) detecting adoption using only config files (Pass A),
versus (2) detecting adoption using the full pipeline (all passes combined). The question is:
does the full pipeline find significantly more agentic repos than Pass A alone?
\end{conceptbox}

\begin{wherebox}
\textbf{Where we use it.}
\begin{itemize}
  \item \textbf{Stage 4b, RQ2 --- Multi-signal vs single-signal:} compares the adoption rate
        found by Pass A alone against the full pipeline on the same corpus.
        Reports the percentage-point difference ($\Delta$pp) and the McNemar $\chi^2$ statistic.
\end{itemize}
\end{wherebox}

%% ═══════════════════════════════════════════════════════════════════════════
\section{Chi-Squared Test ($\chi^2$)}
%% ═══════════════════════════════════════════════════════════════════════════

\begin{conceptbox}
\textbf{What it is.}
Chi-squared tests whether two categorical variables are related or independent. You make a table
of counts (like a spreadsheet with groups in rows and outcomes in columns) and the test asks:
``are the differences between groups too large to be random noise?''

\medskip
\textbf{Simple example.}
You survey people in 5 cities and ask if they own a car (yes/no). Chi-squared tests whether
car ownership rates differ significantly across cities, or whether the differences are just
random fluctuations you would expect from sampling.

\medskip
In the thesis, the ``groups'' are countries (or regions), and the ``outcome'' is agentic
versus non-agentic. Chi-squared is the \emph{omnibus} test --- it tells you that at least one
country differs from the others, but not which one specifically (that is what Fisher does next).
\end{conceptbox}

\begin{wherebox}
\textbf{Where we use it.}
\begin{itemize}
  \item \textbf{Stage 4b, RQ3 --- Geographic variation:} $\chi^2$ across all eligible countries
        to confirm that adoption rates differ significantly by country.
  \item \textbf{Stage 4e (MENA analysis):} within-MENA $\chi^2$ across the 5 sub-regions
        (Gulf, Levant, North Africa, Turkey, Iran).
  \item \textbf{Stage 5g (World regions):} $\chi^2$ omnibus across the 10 macro-regions.
\end{itemize}
\end{wherebox}

%% ═══════════════════════════════════════════════════════════════════════════
\section{Fisher Exact Test}
%% ═══════════════════════════════════════════════════════════════════════════

\begin{conceptbox}
\textbf{What it is.}
Fisher exact is similar to chi-squared but more reliable when sample sizes are small.
Chi-squared is an approximation that becomes inaccurate with small counts (fewer than
5--10 items in a cell). Fisher computes the exact probability, so it works even with
very small groups.

\medskip
Fisher is typically used \emph{pairwise} --- comparing two specific groups against each other,
or one group against a global baseline. It answers: ``is country X's adoption rate
significantly different from the global average?''

\medskip
\textbf{Simple example.}
You have only 12 repos from Luxembourg. Chi-squared would be unreliable here. Fisher exact
handles the small number correctly and tells you whether Luxembourg's rate is genuinely
different from the global rate.
\end{conceptbox}

\begin{wherebox}
\textbf{Where we use it.}
\begin{itemize}
  \item \textbf{Stage 4b, RQ3:} pairwise Fisher for each country vs the global adoption mean,
        after the chi-squared omnibus confirms that variation exists.
  \item \textbf{Stage 4e (MENA):} each MENA country vs the global mean (Fisher exact +
        BH correction). Also used for each MENA sub-region vs global.
  \item \textbf{Stage 5g (World regions):} pairwise Fisher between all pairs of macro-regions
        after the omnibus chi-squared.
\end{itemize}
\end{wherebox}

%% ═══════════════════════════════════════════════════════════════════════════
\section{Benjamini-Hochberg (BH) FDR Correction}
%% ═══════════════════════════════════════════════════════════════════════════

\begin{conceptbox}
\textbf{What it is.}
When you run many statistical tests at the same time --- say, 50 countries, each compared
against the global mean --- you expect roughly 5\% of them to appear significant purely by
chance (since we use a 5\% threshold). That means about 2--3 countries would look significant
even if nothing real is happening. This is called the \emph{multiple testing problem}.

\medskip
Benjamini-Hochberg (BH) fixes this by adjusting the significance thresholds across all your
tests together, controlling the \emph{False Discovery Rate} (FDR) --- the proportion of your
``significant'' results that are actually false alarms.

\medskip
\textbf{Simple example.}
Imagine testing 100 medicines for a side effect using p $<$ 0.05. By pure chance, about 5
medicines will appear to have a side effect even if none of them do. BH correction raises the
bar for the weakest results so that those 5 false alarms are filtered out.

\medskip
\textbf{BH vs Bonferroni:} Another common correction is Bonferroni, which is stricter.
We use BH because it is less conservative --- it allows more true discoveries through while
still controlling false alarms. Bonferroni can miss real effects when you have many tests.
\end{conceptbox}

\begin{wherebox}
\textbf{Where we use it.}
\begin{itemize}
  \item \textbf{Stage 4b:} applied to all pairwise Fisher tests across countries (RQ3).
  \item \textbf{Stage 4e (MENA):} applied to Fisher tests for each MENA country vs global mean.
  \item \textbf{Stage 5g (World regions):} applied to all pairwise Fisher tests across
        the 10 macro-regions.
\end{itemize}
BH correction is always applied together with Fisher exact whenever more than one pairwise
comparison is made.
\end{wherebox}

%% ═══════════════════════════════════════════════════════════════════════════
\section{Kruskal-Wallis Test}
%% ═══════════════════════════════════════════════════════════════════════════

\begin{conceptbox}
\textbf{What it is.}
Kruskal-Wallis tests whether three or more groups have different distributions.
It is the non-parametric version of a one-way ANOVA, meaning it does \emph{not} assume
the data follows a normal (bell-curve) distribution. This matters because repo sizes
(commits, stars, lines of code) are heavily skewed --- a few enormous repos pull the
average far from the typical repo.

\medskip
Kruskal-Wallis works by ranking all values together (smallest = rank 1, largest = rank N)
and then checking whether the ranks are distributed evenly across groups.

\medskip
\textbf{Simple example.}
You have three groups of students: beginners, intermediates, and experts. You measure their
test scores. The scores are not normally distributed (a few geniuses skew the data). Kruskal-Wallis
tests whether at least one group scores significantly differently from the others, without
making assumptions about the shape of the distribution.

\medskip
In the thesis, the three groups are: Small repos / Medium repos / Large repos (split by percentile).
The question is: do larger repos have significantly higher adoption rates?
\end{conceptbox}

\begin{wherebox}
\textbf{Where we use it.}
\begin{itemize}
  \item \textbf{Stage 4d (Size analysis):} tests whether adoption rates differ across three
        repo size strata (small / medium / large), applied separately for three size metrics:
        star count, commit count, and code volume.
  \item \textbf{Stage 5b (Repo size comparison):} 3-way Kruskal-Wallis comparing
        No-AI / AI-Assisted / Agentic repos on four metrics: commits, code KB, stars, forks.
\end{itemize}
\end{wherebox}

%% ═══════════════════════════════════════════════════════════════════════════
\section{Mann-Whitney U Test}
%% ═══════════════════════════════════════════════════════════════════════════

\begin{conceptbox}
\textbf{What it is.}
Mann-Whitney U is a pairwise version of Kruskal-Wallis --- it compares exactly two groups.
You use it after Kruskal-Wallis tells you ``at least one group differs'' to zoom in and
identify which specific pair drives the difference.

\medskip
Like Kruskal-Wallis, it is rank-based and makes no assumptions about the distribution shape.

\medskip
\textbf{Simple example.}
After Kruskal-Wallis tells you that beginners, intermediates, and experts score differently,
Mann-Whitney U lets you specifically compare beginners vs experts to quantify that difference.
\end{conceptbox}

\begin{wherebox}
\textbf{Where we use it.}
\begin{itemize}
  \item \textbf{Stage 5b (Repo size comparison):} pairwise comparison between the
        \emph{Agentic} group and the \emph{No-AI} group specifically, for each of the four
        size metrics (commits, code KB, stars, forks). This is the most informative comparison
        since we want to know how agentic repos differ from repos with no AI at all.
\end{itemize}
\end{wherebox}

%% ═══════════════════════════════════════════════════════════════════════════
\section{Rank-Biserial Effect Size ($r_{rb}$)}
%% ═══════════════════════════════════════════════════════════════════════════

\begin{conceptbox}
\textbf{What it is.}
Statistical significance tells you whether a difference is real (not by chance).
Effect size tells you whether the difference is \emph{meaningful in practice}.
With tens of thousands of repos, even a difference of 1 commit between groups can be
``statistically significant'' --- but that does not mean it matters.

\medskip
Rank-biserial ($r_{rb}$) is an effect size measure for Mann-Whitney U. It ranges from
$-1$ to $+1$:
\begin{itemize}
  \item $|r_{rb}| \approx 0.1$ --- small effect (difference exists but is minor)
  \item $|r_{rb}| \approx 0.3$ --- medium effect (noticeable in practice)
  \item $|r_{rb}| \approx 0.5$ or higher --- large effect (practically meaningful)
\end{itemize}

\medskip
\textbf{Simple example.}
You find that agentic repos have significantly more commits (p $<$ 0.001). But $r_{rb} = 0.08$
means the effect is tiny --- statistically real but practically irrelevant. If $r_{rb} = 0.52$,
the difference is both statistically real and large enough to matter.

\medskip
\textbf{Why it matters for this thesis:}
With $\sim$50{,}000 repos, almost any difference will be statistically significant.
Reporting $r_{rb}$ alongside the p-value prevents over-interpreting trivial differences.
\end{conceptbox}

\begin{wherebox}
\textbf{Where we use it.}
\begin{itemize}
  \item \textbf{Stage 5b (Repo size comparison):} reported alongside Mann-Whitney U for each
        of the four size metrics (commits, code KB, stars, forks) comparing Agentic vs No-AI repos.
        Tells us not just \emph{that} the groups differ in size, but \emph{how much} they differ.
\end{itemize}
\end{wherebox}

%% ═══════════════════════════════════════════════════════════════════════════
\section*{Quick Reference Summary}
%% ═══════════════════════════════════════════════════════════════════════════

\vspace{4pt}
\begin{tabular}{L{3.8cm}L{4cm}L{6.5cm}}
\toprule
\textbf{Test} & \textbf{Question it answers} & \textbf{Used in} \\
\midrule
\rowcolor{rowgray}
Wilson 95\% CI & How confident are we in this rate? & Stages 4a, 4b, 5g \\
Binomial test & Is global adoption $>$ 5\%? & Stage 4b, RQ1 \\
\rowcolor{rowgray}
McNemar test & Does the full pipeline find more than Pass A alone? & Stage 4b, RQ2 \\
Chi-squared ($\chi^2$) & Do rates differ across countries/regions? & Stages 4b, 4e, 5g \\
\rowcolor{rowgray}
Fisher exact & Is country X different from the global mean? & Stages 4b, 4e, 5g \\
BH FDR correction & Are our many Fisher tests producing false alarms? & Stages 4b, 4e, 5g \\
\rowcolor{rowgray}
Kruskal-Wallis & Do small/medium/large repos differ in adoption? & Stages 4d, 5b \\
Mann-Whitney U & Are agentic repos bigger than no-AI repos? & Stage 5b \\
\rowcolor{rowgray}
Rank-biserial $r_{rb}$ & How large is the agentic vs no-AI size difference? & Stage 5b \\
\bottomrule
\end{tabular}

\end{document}
